\section{Compute-Continuum Orchestration and Adaptive Execution}
\label{sec:orchestration}

The optimisation model specifies a placement decision, but a Compute-Continuum workflow also needs an operational procedure for gathering state, selecting a scheduler, executing tasks, and reacting to change. This section defines that procedure at the level of workflow semantics rather than a specific middleware product. Its purpose is to show how the retained formulation and algorithms can be embedded in an eScience orchestration loop that spans edge, fog, cloud, and HPC resources.

\subsection{Orchestration Inputs and Monitoring State}

At each scheduling point, the orchestrator combines a static workflow descriptor with dynamic resource observations. The static descriptor includes the DAG, task identifiers, task dependencies, expected data sizes, execution profiles, eligibility constraints, deadline, and approximation requirements. The dynamic resource state includes the availability of each resource, current queue or load indicator, link bandwidth estimate, link reachability, storage locality, and policy restrictions. A real deployment may obtain these values from a workflow-management service, an instrumentation layer, a resource-information service, or a site-specific monitoring system. The scheduling model does not require a specific source, provided that the values are mapped consistently to the parameters in Section~\ref{sec:system_model}.

The most important monitoring rule is that resource discovery must be coupled with data-path discovery. A newly visible cloud or HPC resource should not automatically be considered a superior candidate merely because it offers higher compute capability. The orchestrator must determine whether a task is allowed to execute there, whether required input data can reach it, and whether the resulting communication delay still fits the deadline. This rule is represented by the eligibility matrix $P_{ia}$ and the data-transfer delay $c_{ji}^{ba}$. It protects the workflow against a common failure mode in heterogeneous systems: moving a task to a powerful remote resource while ignoring the cost of transmitting its prerequisites.

\subsection{Scheduler Selection Policy}

Different workflow states require different optimisation effort. Table~\ref{tab:orchestration_policy} gives an implementation-neutral policy for selecting among the retained methods. The table is not a new performance claim; it translates the strengths and limitations demonstrated by the supplied comparisons into operational guidance. An orchestration service can choose a method before each scheduling run, record the choice in the provenance log, and reproduce the decision later from the same state information.

\begin{table*}[t]
\centering
\caption{Compute-Continuum scheduler-selection policy derived from the retained formulation and reported comparisons.}
\label{tab:orchestration_policy}
\small
\begin{tabular}{p{0.19\textwidth}p{0.25\textwidth}p{0.26\textwidth}p{0.20\textwidth}}
\toprule
\textbf{Operational condition} & \textbf{Preferred method} & \textbf{Rationale} & \textbf{Provenance fields to record} \\
\midrule
Small or stable workflow; sufficient planning time & MILP & Provides a global optimisation benchmark for the current resource and workflow state. & Solver version, model version, time limit, termination status, objective value. \\
\addlinespace
Large workflow; moderate planning time; quality-sensitive result & Randomized Relaxation & Preserves model structure while reducing the difficulty of integral assignment. & Relaxation version, random seed, fractional solution summary, rounding policy. \\
\addlinespace
Tight scheduling deadline or fast local reaction & Greedy scheduler & Makes a low-overhead decision using priority, eligibility, execution time, data size, and bandwidth. & Priority weights, approximation factor, ready-set order, selected placements. \\
\addlinespace
Large and dynamic search space; planning time available & Genetic-Algorithm Metaheuristic & Explores alternative assignments and approximation levels beyond local choices. & Population size, generation count, crossover/mutation settings, seed, fitness trace. \\
\addlinespace
Node failure, bandwidth degradation, or policy change & Re-evaluate eligible methods on updated state & Removes unusable resources and recomputes data-path-aware placement under the remaining deadline. & Failure event, updated resource graph, discarded placements, recovery schedule. \\
\bottomrule
\end{tabular}
\end{table*}

The policy makes explicit that no single method is optimal for every orchestration event. The MILP is valuable when the decision interval permits a detailed global calculation. The randomized relaxation can be selected when the workflow is larger, but the orchestrator still requires a model-informed placement. Greedy scheduling is appropriate for a rapid local response, and the genetic algorithm is appropriate when the search space is broad enough to justify its iterative exploration. Recording this choice is essential: otherwise, two nominally identical workflow executions may be impossible to compare because the scheduler or configuration differed.

\subsection{Adaptive Execution and Resilience Loop}

A scheduled workflow should be monitored throughout execution rather than treated as a one-time placement. A resource failure, a significant decrease in bandwidth, an increase in data volume, or a deadline revision can invalidate the assumptions that produced the original schedule. The orchestrator therefore follows a four-step loop. First, it detects a state change and identifies the affected tasks and data dependencies. Second, it updates resource eligibility and communication estimates. Third, it determines whether the remaining portion of the workflow can continue under the existing schedule or requires replanning. Fourth, it invokes the selected scheduling method for the unresolved tasks and records the reconfiguration decision.

This loop preserves the non-preemptive assumption in the original formulation. A task that has started is not silently migrated; instead, the scheduler replans tasks that have not started or that have not yet received all required predecessor data. If an application permits checkpoint-and-restart, a checkpoint can be represented as a new data-producing task in the DAG. This makes restart cost and data movement visible to the same scheduling logic instead of treating recovery as an external, unmodelled action.

The reported failure-rate comparison in Figure~\ref{fig:perf_failure_accuracy} motivates this loop. It shows the loss of reported average accuracy as the failure rate increases, and it establishes a basis for selecting higher-quality scalable methods when the failure pressure becomes material. It does not measure the detection, restart, or state-replication time of a production system. Those quantities should be added to a future trace-level deployment study.

\subsection{Data Streaming and In-Transit Processing}

The DAG representation also supports a streaming interpretation. A streaming source can emit a sequence of data objects that are treated as repeated workflow instances or as windows in a larger workflow. For each window, the scheduler can update $D_{ji}$ from the observed data volume and update $B_{ba}$ from the current network measurements. A task may then be placed near the source for immediate processing, forwarded to a federated fog node for collaborative analysis, or sent to cloud/HPC capacity when its latency budget permits. In-transit processing is represented by inserting transformation tasks on the data path, thereby making their compute cost and reduced output size explicit.

This representation is useful because it prevents a data-streaming decision from becoming an opaque networking choice. A compression, filtering, aggregation, or feature-extrac\-tion task has a measurable execution time and changes the data size passed to downstream tasks. The scheduler can consequently evaluate whether performing that task near the source reduces the total deadline risk. The same workflow metadata records the transformation, which improves the interpretability and reproducibility of the data path.

\subsection{Example End-to-End Execution Trace}

Consider a disaster-response workflow that begins with distributed sensor observations and ends with a decision-support output. The edge tier first records the observation time, source identifier, data schema, and data-volume estimate. A lightweight preprocessing task can validate the input and remove malformed records. The resulting window is then registered as a DAG instance with a deadline associated with the operational response requirement. The registration event is a provenance boundary: it identifies the exact input window and the workflow version before scheduling begins.

The orchestrator next queries the current fog federation. It receives the availability and eligibility of each candidate node, the estimated bandwidth of links that might carry intermediate data, and any site-level policy constraints. The scheduler then creates a placement plan. A feature-extraction task may remain close to the source because forwarding raw data would create a large communication delay. A downstream simulation task may be placed on a more capable fog resource when the reduced feature representation makes the transfer feasible. A nonurgent aggregation task may be assigned to a cloud or HPC endpoint if the workflow's remaining deadline accommodates the data path. Each of these choices is stored together with the parameter values that justified it.

During execution, a monitoring event can cause the plan to change. If a fog node fails before a dependent task begins, the orchestrator updates the resource graph, removes the failed node from eligible placements, recalculates the relevant transfer delays, and invokes the selected scheduler for the remaining DAG. The updated plan preserves the identifiers of completed tasks and the data products already available, so re-execution is limited to unresolved work. If a task has already started, the non-preemptive assumption keeps it in place; a future checkpoint-and-restart extension can expose a checkpoint as a new task in the DAG.

At completion, the workflow produces a run manifest that contains the input identifier, DAG descriptor version, resource-state snapshot, selected method, scheduler configuration, random seed where applicable, task placements, completion times, and result metrics. The manifest links to the plots or tables generated from the run. This trace permits a later evaluator to distinguish a difference caused by workflow input, resource state, policy, or scheduler version. It also supports the use of the workflow as a reusable research object: a new site can preserve the same DAG and configuration while substituting an appropriate resource descriptor for its own Compute-Continuum deployment.

\subsection{Energy, Cost, and Security Policy Hooks}

The supplied results evaluate accuracy, scale, and resilience rather than measured energy consumption or monetary cost. The model should therefore not be presented as an energy- or cost-optimal method without additional data. Nevertheless, an implementation can extend the resource descriptor with energy intensity, power cap, carbon signal, price, data-egress charge, and security classification. These attributes can enter the eligibility policy, a secondary objective, or a constraint in a future extension. For example, a task with regulated data may have an eligibility constraint that excludes resources outside an approved trust boundary; a cost-sensitive batch stage may use a maximum data-egress budget; and an energy-sensitive workflow may prefer resources with a stated power policy.

By treating these factors as explicit metadata rather than informal deployment assumptions, the workflow remains interoperable across sites, and its policy choices remain auditable. A future experimental study should collect the necessary resource traces and report the trade-offs between accuracy, makespan, energy, cost, and security constraints. This is a natural extension of the provenance methodology in Section~\ref{sec:reproducibility}.
